
%%%%%%%%%%%%%%%%%%%%%%%%%%%%%%%%%%%%%%%%
%           main body
%%%%%%%%%%%%%%%%%%%%%%%%%%%%%%%%%%%%%%%%

%-----------------------------------
%	TITLE PAGE
%-----------------------------------

\title[数分II\\
复习]{\texorpdfstring{\Huge \bfseries 数学分析II复习}{数学分析II复习}} % The short title appears at the bottom of every slide, the full title is only on the title page
% \author[\S 11.3\\
% 多元连续函数性质] % Your institution as it will appear on the bottom of every slide, maybe shorthand to save space
% {\Large \bfseries \S 11.3 多元连续函数的性质}
% \author{Singular Value Decomposition} % Your name
% \date{\today} % Date, can be changed to a custom date
\date{}

%------------------------------------------------

\begin{document}

%------------------------------------------------

\begin{frame}

\titlepage % Print the title page as the first slide

\end{frame}

%------------------------------------------------

% \begin{frame}
% \frametitle{内容提要} % Table of contents slide, comment this block out to remove it
% \tableofcontents % Throughout your presentation, if you choose to use \section{} and \subsection{} commands, these will automatically be printed on this slide as an overview of your presentation
% \end{frame}

%-----------------------------------
%	PRESENTATION SLIDES
%-----------------------------------

\section{定积分}

%------------------------------------------------

\begin{frame}{可积性与积分的定义}

基本考察对象: 定义在闭区间上函数 $f: [a, b] \longrightarrow \mathbb{R}$
\pause
$\longsquiggly$ $\mathcal{R}[a,b].$

\pause

\begin{itemize}[<+->]
\item {\color{red}带标志点的}划分 $(P, {\color{red}\xi}):$
\begin{alignat*}{8}
P & : a = x_0 & < & ~ x_1 & < & ~ x_2 & ~ < & ~ \cdots & ~ < & ~ x_n = b \\
& & \uparrow & & \uparrow & & \uparrow & & \uparrow & \\
{\color{red}\xi} & : & ~ {\color{red}\xi_1} & & ~ {\color{red}\xi_2} & & ~ {\color{red}\xi_3} & ~ \cdots & ~ {\color{red}\xi_n} &
\end{alignat*}
$\Delta_i = [x_{i-1}, x_i], ~ \Delta x_i = x_i - x_{i-1}, ~ \lambda(P) = \max_{1 \leqslant i \leqslant n} \Delta x_i$
\item 积分和: $S(f; P; \xi) = \sum\limits_{i=1}^n f(\xi) \Delta x_i.$
\item Darboux 大和: $\overline{S}(f; P) = \sum\limits_{i=1}^n \sup_{t\in\Delta_i} f(t) \cdot \Delta x_i =: \sum\limits_{i=1}^n M_i(f) \cdot \Delta x_i$ \par
Darboux 小和: $\underline{S}(f; P) = \sum\limits_{i=1}^n \inf_{t\in\Delta_i} f(t) \cdot \Delta x_i =: \sum\limits_{i=1}^n m_i(f) \cdot \Delta x_i$
\end{itemize}

\end{frame}

%------------------------------------------------

\begin{frame}{可积性与积分的定义}

\begin{Def}[函数 Riemann 可积]
极限 $\lim\limits_{\lambda(P) \to 0} S(f; P; \xi) = I \in \mathbb{R}$ 存在, 则称 $f$ 在 $[a, b]$ 上 Riemann 可积, 记作
$$\int_a^b f(x) \mathrm{d} x = \lim\limits_{\lambda(P) \to 0} S(f; P; \xi).$$
\end{Def}

\end{frame}

%------------------------------------------------

\begin{frame}{可积函数的性质与判别}

必要条件: $f(x)$ 有界.

\pause
\vspace{1em}

等价 (充要) 条件:
\begin{enumerate}[<+->]
\item Cauchy 收敛原理: $\forall \varepsilon > 0, \exists \delta > 0,$ 使得任意带标志点的划分 $(P_1, \xi_1), (P_2, \xi_2),$ 若满足 $\lambda(P_1) < \delta, \lambda(P_2) < \delta,$ 总有
$$\lvert S(f; P_1; \xi_1) - S(f; P_2; \xi_2) \rvert < \varepsilon.$$
\item Darboux 和: $\lim\limits_{\lambda(P) \to 0} \overline{S}(f; P) =: \overline{\int}_a^b f(x) \mathrm{d} x = \underline{\int}_a^b f(x) \mathrm{d} x := \lim\limits_{\lambda(P) \to 0} \underline{S}(f; P).$
\item 振幅: $\forall \varepsilon > 0, \exists \delta > 0,$ 使得任意满足 $\lambda(P) < \delta$ 的划分 $P$ 总有
$\sum\limits_{i=1}^n \omega(f; \Delta_i) \Delta x_i < \varepsilon,$ 振幅 $\omega(f; \Delta_i) := \sup\limits_{t_1, t_2 \in \Delta_i} \lvert f(t_1) - f(t_2) \rvert.$
\item Lebesgue 判别法: 几乎处处连续 (不连续点集是零测集).
\end{enumerate}

\end{frame}

%------------------------------------------------

\begin{frame}{闭区间上典型的可积函数}

\begin{itemize}[<+->]
\item 连续函数 (所有初等函数)
\item 单调函数
\item $[0, 1]$ 上的 Riemann 函数 $R(x) = \begin{cases}
  \dfrac{1}{q}, & x = \dfrac{p}{q} ~ \text{为既约分数;}\\
  0, & x ~ \text{为无理数.}
\end{cases}$
\end{itemize}

\pause
\vspace{1em}

典型的不可积函数:

\begin{itemize}
\item $[0, 1]$ 上的 Dirichlet 函数 $D(x) = \begin{cases}
  1, & x ~ \text{为有理数;}\\
  0, & x ~ \text{为无理数.}
\end{cases}$
\item 无界函数, 例如 $f(x) = \begin{cases}
  \dfrac{1}{x}, & x \in (0, 1];\\
  0, & x = 0.
\end{cases}$
\end{itemize}

\end{frame}

%------------------------------------------------

\begin{frame}{定积分的 (运算) 性质}

关于被积区间的
\pause
\begin{enumerate}[<+->]
\item 区间 (有限) 可积性: $\int_a^b f(x) \mathrm{d} x = \int_a^c f(x) \mathrm{d} x + \int_c^b f(x) \mathrm{d} x.$
\end{enumerate}

\pause

关于被积函数的
\pause
\begin{enumerate}[<+->]
\item 线性性:
$$\int_a^b (\alpha g(x) + \beta f(x)) \mathrm{d} x = \alpha \int_a^b g(x)  \mathrm{d} x + \beta \int_a^b f(x) \mathrm{d} x.$$
\item 乘积可积性: $f, g \in \mathcal{R}[a, b] \Longrightarrow f \cdot g \in \mathcal{R}[a, b].$
\item 保序性: $f \leqslant g \Longrightarrow \int_a^b f(x) \mathrm{d} x \leqslant \int_a^b g(x) \mathrm{d} x.$
\item 绝对可积性: $f \in \mathcal{R}[a, b] \Longrightarrow \lvert f \rvert \in \mathcal{R}[a, b]$ (反之不成立), 且有
$$\left\lvert \int_a^b f(x) \mathrm{d} x \right\rvert \leqslant \int_a^b \lvert f(x) \rvert \mathrm{d} x.$$
\end{enumerate}

\end{frame}

%------------------------------------------------

\begin{frame}{积分中值定理}

\begin{enumerate}[<+->]
\item 积分第一中值定理: $f, g \in \mathcal{R}[a, b],$ $g$ 不变号, 那么存在 $\inf_{x \in [a,b]} f(x) \leqslant \eta \sup_{x \in [a,b]} f(x),$ 使得 $\int_a^b f(x) g(x) \mathrm{d} x = \eta \int_a^b g(x) \mathrm{d} x.$
\item 积分第二中值定理
\end{enumerate}

\end{frame}

%------------------------------------------------

\begin{frame}{定积分的计算}

\end{frame}

%------------------------------------------------

\begin{frame}{定积分的应用}

\end{frame}

%------------------------------------------------

\section{反常积分}

%------------------------------------------------

\begin{frame}{敛散性判别法}


\end{frame}

%------------------------------------------------

\begin{frame}{待写}


\end{frame}

%------------------------------------------------

\section{数项级数}

%------------------------------------------------

\begin{frame}{正项级数敛散性判别法}



\end{frame}

%------------------------------------------------

\begin{frame}{一般级数敛散性判别法}


\end{frame}

%------------------------------------------------

\begin{frame}{无穷乘积}

\end{frame}

%------------------------------------------------

\section{函数项级数}

%------------------------------------------------

\begin{frame}{一致收敛性及判别法}



\end{frame}

%------------------------------------------------

\begin{frame}{一致收敛级数的分析性质}



\end{frame}

%------------------------------------------------

\begin{frame}{幂级数}



\end{frame}

%------------------------------------------------

\section{Euclid 空间}

%------------------------------------------------

\begin{frame}{Euclid 空间结构}



\end{frame}

%------------------------------------------------

\begin{frame}{紧集上的连续函数}

\end{frame}

%------------------------------------------------

\begin{frame}{道路连通性}



\end{frame}

%------------------------------------------------

\end{document}
